% ===========================================================================
\section{熵：平均「surprise」}\label{sec:entropy}
\warmth{0}\quad\whisper{熵把所有可能结果的「surprise」取平均。}

\begin{strand}
$-\log p(x)$ 描述某一个结果的「surprise」。熵对它取期望，因此熵依赖于
$X$ 的\emph{分布}，而不依赖于这一次实际观察到了哪个取值。
\end{strand}

\begin{definition}[$b$ 进制熵]\label{def:entropy}
设离散随机变量 $X$ 的字母表为 $\X$，则
\[
  \Ent_b(X)
  =-\sum_{x\in\X}p_X(x)\log_b p_X(x)
  =\E\!\left[-\log_b p_X(X)\right].
\]
约定 $0\log_b0=0$。若无特别说明，取 $b=2$。
\end{definition}

\begin{remark}
熵只依赖于 $p_X$，因此 $\Ent(X)$、$\Ent(p_X)$ 和 $\Ent(p)$
可能表示同一个数。
\end{remark}

\block{伯努利随机变量}

\begin{example}[二元熵]
设 $\X=\{H,T\}$，并且 $\PP(X=H)=p$、$\PP(X=T)=1-p$。则
\[
  \Ent(X)=-p\log p-(1-p)\log(1-p).
\]
在这个特殊情形下，也常把它写成 $\Ent(p)$。
\end{example}

\begin{yourturn}
用公式补全三个关键取值：
\[
  \Ent(0)=\answerin[1cm]{0},\qquad
  \Ent\!\left(\tfrac12\right)=\answerin[1cm]{1},\qquad
  \Ent(1)=\answerin[1cm]{0}\quad\text{比特}.
\]
画出 $0\le p\le1$ 时的大致曲线，并标出最高点。
\workspace[3]
\recall{三个值依次是 $0,1,0$。曲线关于 $p=1/2$ 对称，并在这里达到最高点。}
\end{yourturn}

\block{联合熵}

\begin{definition}[联合熵]\label{def:joint-entropy}
令 $X=(X_1,X_2)$，其联合概率质量函数为 $p_{X_1X_2}$，则
\[
  \Ent(X_1,X_2)
  =-\sum_{x_1\in\X_1}\sum_{x_2\in\X_2}
    p_{X_1X_2}(x_1,x_2)\log p_{X_1X_2}(x_1,x_2).
\]
\end{definition}

\begin{proposition}[独立时的可加性]\label{prop:independent-entropy}
若 $X_1$ 与 $X_2$ 独立，则
\[
  \Ent(X_1,X_2)=\Ent(X_1)+\Ent(X_2).
\]
\end{proposition}

\begin{proof}
独立性给出
$p_{X_1X_2}(x_1,x_2)=p_{X_1}(x_1)p_{X_2}(x_2)$。
把它代入定义~\ref{def:joint-entropy}，并拆开乘积的对数：
\begin{align*}
  \Ent(X_1,X_2)
  &=-\sum_{x_1,x_2}p_{X_1}(x_1)p_{X_2}(x_2)
    \bigl(\log p_{X_1}(x_1)+\log p_{X_2}(x_2)\bigr)\\
  &=-\sum_{x_1}p_{X_1}(x_1)\log p_{X_1}(x_1)
      \sum_{x_2}p_{X_2}(x_2)\\
  &\quad-\sum_{x_2}p_{X_2}(x_2)\log p_{X_2}(x_2)
      \sum_{x_1}p_{X_1}(x_1)\\
  &=\Ent(X_1)+\Ent(X_2).
\end{align*}
最后一步使用了两个边缘概率质量函数之和均为 $1$。
\end{proof}

\begin{yourturn}
复述可加性证明中最关键的一步：
\[
  \log\!\bigl(p_{X_1}(x_1)p_{X_2}(x_2)\bigr)
  =\answerin[5.2cm]{\log p_{X_1}(x_1)+\log p_{X_2}(x_2)}.
\]
\recall{$\log p_{X_1}(x_1)+\log p_{X_2}(x_2)$。}
\end{yourturn}
